\section{Geometry and constructions of the two laws}\label{sup:geometry}

% The paper cites these subsections by number (S1.1 to S1.14): keep their order.

This section gives the constructions that Section~3 of the paper names: how the overlap is
computed, how bodies made of convex pieces are treated, how the pressure of the distributed law
is integrated exactly, and the details of the friction, of the stiction and of the baselines.
Table~\ref{tab:notation-full} lists the complete notation.

\begin{table}[!htbp]
\caption{Notation. Vectors are bold; all quantities are expressed in the world frame unless a
frame is named}\label{tab:notation-full}
\footnotesize
\begin{tabular}{@{}lp{0.70\linewidth}@{}}
\toprule
Symbol & Meaning \\
\midrule
$A$, $B$ & the two bodies of a pair (or two of their convex pieces) \\
$A_c$, $A_T$ & area of a flat contact; area of a triangle \\
$\region$, $\partial\region$ & overlap $A\cap B$ of the undeformed bodies; its boundary \\
$\partial\region_A$, $\partial\region_B$ & parts of $\partial\region$ on the surface of $A$ (inside $B$) and of $B$ (inside $A$) \\
$V$, $\vect{c}$, $\tens{J}$ & volume, centroid and inertia tensor (unit density, about $\vect{c}$) of $\region$ \\
$\n$ & pair normal, unit vector that pushes $A$ out of $B$ \\
$\vect{S}_B$ & vector area of $\partial\region_B$ \\
$\vv(\x)$, $\w$ & velocity at $\x$ and angular velocity of $A$ relative to $B$ \\
$v_n$, $\vv_t$ & $\n\cdot\vv$ (positive when separating) and the tangential part of $\vv$ \\
$\omega_n$, $\w_t$ & $\n\cdot\w$ and the tangential part of $\w$ \\
$K$, $h$ & stiffness [N/m$^3$] and damping factor [s/m] of the foundation \\
$\mu$, $\mu_s$ & kinetic and static coefficients of friction \\
$\ell(\x)$ & length of the column of $\region$ through $\x$ along $\n$ \\
$\varphi(\x)$ & damping factor of the pressure, $1 - h\,\n\cdot\vv(\x)$ \\
$\gamma$ & share of a column carried by one of its ends, $\tfrac12$ or 1 \\
$p(\x)$, $\Fn$ & pressure; normal force \\
$\m$ & unit normal of a face or triangle \\
$\q$, $\rr$, $\eta$ & instantaneous centre of the slip in the plane of a triangle; $\rr = \x - \q$; distance from $\q$ to the line of an edge \\
$D$, $\vect{R}$, $\vect{U}$, $\tens{W}$ & polar moments of a triangle about $\q$ \\
$r_g$ & radius of gyration of $\region$ about $\n$, $r_g^2 = \n\cdot\tens{J}\,\n / V$ \\
$\rho$, $\vect{c}_p$ & radius of gyration of the pressure about its centre; centre of the normal load \\
$\vect{M}_r$, $\vect{M}_s$ & rolling-resistance and spinning-friction torques \\
$d$, $\tau$ & pre-sliding displacement at the static limit; time constant of the stick damper \\
$v_s$, $s$ & stick--slip transition velocity; weight of the stick force, between 0 (sliding) and 1 (stuck) \\
$\dt$, $w$ & time step; inverse mass seen by a contact [1/kg] \\
$\Lambda$ & load acceleration of a body, $\sum|\Fn|\,w$ \\
\botrule
\end{tabular}
\end{table}

\subsection{Shapes, acceleration structures and clipping}

A shape is a closed, manifold triangle mesh with outward orientation, given in the frame of its
body. Three structures are built from it once. The first is a tree of axis-aligned bounding
boxes over its triangles \cite{vandenBergen1997}, stored depth-first in one array. The second
is a binary space partition (BSP) of the solid \cite{Thibault1987} whose splitting planes are
the planes of the faces and whose leaves are either inside or outside the solid; it answers
whether a point is inside and cuts a convex polygon into the fragments that lie inside. Its
planes have a thickness of $10^{-10}$ times the diagonal of the bounding box of the shape:
points within that distance of a plane are on it. The third exists only when the mesh is convex
and is the list of its distinct face planes.

The computation follows the 2018 thesis \cite{Vilela2018} in its structure. The region where
the overlap can be is the intersection of the two bounding boxes, each expressed in the frame
of the other and intersected again. The triangles of $A$ whose boxes meet that region are
transformed to the frame of $B$ and clipped by the BSP of $B$, plane by plane, with the
re-entrant polygon clipping of Sutherland and Hodgman \cite{Sutherland1974}; every fragment
that reaches an inside leaf is a convex polygon of $\partial\region_A$. The triangles of $B$
are clipped against $A$ in the same way and give $\partial\region_B$. Each polygon keeps the
outward normal of its face and whether it belongs to $A$ or to $B$.

The BSP of a convex mesh degenerates into a chain: every face is a node whose front subtree is
empty, and a polygon inside the solid walks all of it. For convex shapes with more than 200
faces the polygon is instead cut by the planes of the faces near it, found with the box tree,
followed by one test of whether what is left lies inside the shape. The result is the same as
that of the BSP, which the tests of the library assert.

\subsection{Coincident faces}

Two solids that touch along a face need a rule, because a polygon lying on a splitting plane is
neither in front of it nor behind it, and the patch must contribute to $\partial\region$ once
or not at all (the consistency problem discussed in the 2018 thesis \cite{Vilela2018}). A polygon of
$A$ on a plane of $B$ counts as inside where $B$ is solid just behind it, so a face of $A$ that
coincides with a face of $B$ of the same orientation is kept; a polygon of $B$ on a plane of
$A$ counts only where $A$ is solid on both sides, so a face of $B$ that coincides with any face
of $A$ is dropped. With one plane thickness for both passes (the larger of the two shapes'),
the two passes agree on what is coplanar, and $\partial\region$ is closed up to round-off.

\subsection{Volume integrals and the pair normal}

The polygons are not stored for the volume integrals. As each one is produced it is split into
a fan of triangles, and each triangle adds its term to the integrals of $1$, $x$, $y$, $z$ and of
the six quadratic monomials over $\region$ \cite{Mirtich1996,Eberly2002}. The sums are
accumulated about a point inside the region, which keeps the second moments well conditioned
for overlaps that are thin slabs a few micrometres thick.

For the direction, Gonthier's thesis takes the eigenvector of $\tens{J}$ nearest to a geometric
estimate \cite{Gonthier2007thesis}, and the 2018 thesis fitted a plane to the contour where the
two surfaces intersect \cite{Vilela2018}. Here it is the direction of the vector area
$\vect{S}_B$ of $\partial\region_B$ (Eq.~1 of the paper). Since $\partial\region$ is closed,
the vector areas of its two parts are opposite, $\vect{S}_A = -\vect{S}_B$, and the
implementation uses $(\vect{S}_B - \vect{S}_A)/2$, in which the round-off of the two passes
averages. The vector area of a surface depends only on its boundary curve: when the contour is
planar, $\n$ is the normal of its plane, which is what the least-squares plane of the thesis
approximated, obtained here without extracting the contour; for an edge or a vertex of one body
pressed into a face of the other it is a weighted mean of the normals involved. Each polygon
returns its vector area as a by-product of its volume terms.

\subsection{Bodies made of convex pieces: seams}

A non-convex body is given as convex pieces that fill it without overlapping. The pieces of the
battery are generated procedurally and are exact. An approximate convex decomposition of a
general mesh \cite{Wei2022coacd} gives pieces that overlap or leave gaps; the treatment of this
section does not cover them, and the laws were not tested on them. One pair is evaluated for
every combination of a piece of one body with a piece of the other, after a broad phase that
keeps only the pieces whose boxes are near.

Cutting a body introduces faces that are not surface of the body: the \emph{seams} where two
sibling pieces meet (Fig.~\ref{fig:seams}). The seams of a piece $B_i$ that lie inside $A$ are
part of the boundary of the overlap $A \cap B_i$, and they must stay in its volume integrals,
which need a closed surface. They are not contact surface. When a compound body is built, every
face of every piece is clipped against the sibling pieces with a third coplanar rule (a polygon
counts as inside where the sibling is solid just in front of it); the fragments that have a
sibling in front are flagged as seams, and faces that are seams only in part are
re-triangulated. The flag travels with every polygon of $\partial\region$. Two quantities are
then defined without the seams. The first is the pair normal used by the distributed law, the
direction of the vector area of the faces of $B$ in $\partial\region$ that are surface of the
body: for a piece of a faceted ground it is the normal of the facets in contact, whatever part of
the piece the overlap covers, whereas the vector area with the seams tilts with the depth of the
overlap at the edge of a piece. The second is the pressure: seams carry none and end no column
(Section~\ref{sup:columns}). Gonthier's law, which uses only $V$, $\vect{c}$, $\tens{J}$ and a
normal, keeps the full vector area.

\begin{figure}[!htbp]
\centering
% A body on a ground cut into two convex pieces: seam, columns and normals.
\begin{tikzpicture}[scale=1.0, font=\footnotesize, >=Stealth]
  % Ground pieces B1 and B2.
  \fill[gray!18] (-4,-1.5) rectangle (0.6,0);
  \fill[gray!28] (0.6,-1.5) rectangle (4,0);
  \draw[thick] (-4,0) -- (4,0);
  \draw (-4,-1.5) -- (-4,0) (4,-1.5) -- (4,0);
  \draw[densely dotted] (0.6,-1.5) -- (0.6,-0.45);
  \node at (-3.5,-1.2) {$B_1$};
  \node at (3.5,-1.2) {$B_2$};
  % Body A, flat bottom pressed in by 0.45.
  \fill[gray!8] (-2.6,-0.45) rectangle (2.4,1.6);
  \draw[thick] (-2.6,-0.45) rectangle (2.4,1.6);
  \node at (2.05,1.3) {$A$};
  % Overlaps with each piece.
  \fill[pattern=north east lines, pattern color=gray!70] (-2.6,-0.45) rectangle (0.6,0);
  \fill[pattern=north west lines, pattern color=gray!70] (0.6,-0.45) rectangle (2.4,0);
  % Seam inside the overlap.
  \draw[very thick, densely dashed] (0.6,-0.45) -- (0.6,0);
  \node[anchor=north] at (0.6,-0.55) {seam};
  % Columns of the whole bodies.
  \foreach \xx in {-2.2,-1.8,-1.4,-1.0,-0.6,-0.2,0.2,1.0,1.4,1.8,2.2} {
    \draw[thin] (\xx,-0.45) -- (\xx,0);
  }
  % Normals: with the seam (tilted) and without (law).
  \draw[->, very thick] (-1.0,0.15) -- ++(0,1.0);
  \draw[->, thick, gray] (-1.0,0.15) -- ++(0.28,0.96);
  \draw[->, very thick] (1.5,0.15) -- ++(0,1.0);
  \draw[->, thick, gray] (1.5,0.15) -- ++(-0.42,0.9);
  \node[anchor=east, inner sep=2pt] at (-1.0,0.85) {$\n$};
  \node[gray, anchor=west, inner sep=2pt] at (-0.62,1.2) {with the seam};
  % Labels.
  \node[anchor=west] at (2.6,-0.25) {$A\cap B_2$};
  \node[anchor=east] at (-2.8,-0.25) {$A\cap B_1$};
\end{tikzpicture}

\caption{A body $A$ on a ground $B$ cut into the convex pieces $B_1$ and $B_2$ (section). The
overlap of $A$ with each piece is closed by a seam face (dashed), which enters the volume
integrals of the piece but carries no pressure and does not end a column: the columns (thin
lines) are those of the whole bodies. With the seam, the vector area of the faces of $B_1$ in
the overlap tilts by the ratio of the seam to the contact face; the normal of the distributed
law (arrows) leaves the seam out}
\label{fig:seams}
\end{figure}

\subsection{Deep overlaps}

The vector area $\vect{S}_B$ cancels when the part of $\partial\region$ on $B$ closes on itself
or is empty: a body entirely inside another, or one that pierces another through. Near those
configurations the normal is ill-conditioned, and a law that dropped the pair there would lose
its force at the deepest penetrations. The intersection therefore reports no normal when
$|\vect{S}_B| < 0.1\, V^{2/3}$ (a compact overlap bounded by $B$ on one side has a ratio of
about 2: 1.9 for a half sphere, 2 for the corner of a cube), and the pair then keeps the normal
of its last evaluation, stored in the frame of $B$ as part of the state of the pair: the
direction along which the body came in. A pair that starts that deep takes the outward normal
of the nearest surface face of a convex $B$. Both laws use this fallback.

\subsection{Columns and the exact partition}\label{sup:columns}

The distributed law requires convex shapes (or convex pieces), so that the overlap is convex
and every column is a single segment. A column through a point $\x$ of the contact surface runs
along $\vect{u} = +\n$ from a face of $A$ and along $\vect{u} = -\n$ from a face of $B$, and
ends at the first face of the surface of either body. A face with outward unit normal $\m_f$
and plane $\m_f\cdot\vect{y} = e_f$ ends it only if $\m_f\cdot\vect{u} > 0$, at the distance
$(e_f - \m_f\cdot\x)/(\m_f\cdot\vect{u})$, which is an affine function of $\x$. The column
length is the smallest of these,
\begin{equation}
  \ell(\x) = \min_{f\,:\;\m_f\cdot\vect{u} > 0}\; \frac{e_f - \m_f\cdot\x}{\m_f\cdot\vect{u}} ,
  \label{eq:column}
\end{equation}
over the faces of both shapes that are surface of their bodies. Seams are left out of the
minimum: behind a seam the body continues, so the columns are those of the whole bodies however
they are cut (Fig.~\ref{fig:seams}). Only the faces that can be the first need to be examined:
those that bound the overlap and, for a shape without seams, those within the reach of the
polygon, found with the box tree.

A column carries $K\ell$. When its two ends lie on different bodies, each end carries half
($\gamma = \tfrac12$), so the law is symmetric in $A$ and $B$. When both ends lie on the same
body, which happens when the overlap encloses a body or pierces it, the end that pushes $A$ out
carries it whole ($\gamma = 1$), so that the force stays $K V$ and is continuous through the
enclosure.

On a polygon of $\partial\region$, $\ell$ is the minimum of a few affine functions. The polygon
is split into the convex regions where each face is the one that ends the column, by clipping
it against the lines where two of the affine functions are equal (Fig.~\ref{fig:partition}). On
each region $\ell$ is affine, so the integrals of the elastic pressure and of its moment are
those of a linear function over a polygon and have no quadrature error, whatever the size of
the triangles of the mesh.

\begin{figure}[!htbp]
\centering
% A polygon of the contact surface split by the face where its columns end, and cut where g = 0.
\begin{tikzpicture}[scale=1.0, font=\footnotesize, >=Stealth]
  % Polygon (a face fragment), seen along n.
  \coordinate (v1) at (0,0);
  \coordinate (v2) at (5.2,0.2);
  \coordinate (v3) at (5.8,2.6);
  \coordinate (v4) at (2.6,3.6);
  \coordinate (v5) at (-0.4,2.4);
  % Interior point where the three affine functions are equal, and the points where the
  % dividing lines leave the polygon.
  \coordinate (c) at (2.7,1.7);
  \coordinate (e1) at ($(v1)!0.52!(v2)$);
  \coordinate (e2) at ($(v3)!0.45!(v4)$);
  \coordinate (e3) at ($(v5)!0.3!(v1)$);
  \fill[gray!10] (v1) -- (e1) -- (c) -- (e3) -- cycle;
  \fill[gray!25] (e1) -- (v2) -- (v3) -- (e2) -- (c) -- cycle;
  \fill[gray!40] (c) -- (e2) -- (v4) -- (v5) -- (e3) -- cycle;
  \draw[thick] (v1) -- (v2) -- (v3) -- (v4) -- (v5) -- cycle;
  \draw (c) -- (e1) node[below right, inner sep=2pt] {$\ell_1 = \ell_2$};
  \draw (c) -- (e2) node[above, inner sep=2pt] {$\ell_2 = \ell_3$};
  \draw (c) -- (e3) node[left] {$\ell_1 = \ell_3$};
  \node at (1.5,0.9) {$\ell = \ell_1$};
  \node at (4.4,1.3) {$\ell = \ell_2$};
  \node at (1.6,2.75) {$\ell = \ell_3$};
  % Line g = 0 cutting a corner.
  \coordinate (g1) at ($(v1)!0.2!(v2)$);
  \coordinate (g2) at ($(v1)!0.45!(v5)$);
  \draw[densely dashed, thick] ($(g1)!1.25!(g2)$) -- ($(g1)!-0.3!(g2)$) node[below, inner sep=2pt] {$g = 0$};
  \fill[white, opacity=0.75] (g1) -- (v1) -- (g2) -- cycle;
  \draw[thick] (g1) -- (v1) -- (g2);
  \node[anchor=east] at (-0.3,0.3) {$p = 0$};\draw[thin] (-0.3,0.3) -- (0.2,0.35);
\end{tikzpicture}

\caption{A polygon of the contact surface, seen along $\n$. The column length $\ell$ is the
minimum of the affine functions $\ell_1$, $\ell_2$, $\ell_3$ of the three faces where a column can
end; the polygon is split along the lines where two of them are equal (solid lines), and on each
part $\ell$ is affine. Where the damping factor $\varphi$ changes sign the part is cut again (dashed
line) and the tensile side carries no pressure}
\label{fig:partition}
\end{figure}

\subsection{Damping along the pair normal and the exact clamp}

The rate in the damping factor $\varphi = 1 - h\,\n\cdot\vv$ is measured along the pair normal,
not along the normal of the face, for the following reason. The two ends of the column through
$\x$ are material points of $A$ and of $B$, a distance $\ell$ apart along $\n$. The relative
velocity field is rigid, so the velocities at the two ends differ by $\w\times\ell\n$, which has
no component along $\n$: the quantity $\n\cdot\vv$ is the same all along the column, and it is
the rate at which the column shortens. Sliding along $\n$'s tangent plane and spinning about
$\n$ therefore damp nothing, on any face. Measured along the normal of each face instead, the
slip along any face not perpendicular to $\n$ would count as an approach velocity, which on the
seams of a many-piece bearing is a viscous drag larger than the friction itself.

Without adhesion the pressure cannot be negative. The factor $\varphi$ is affine on the whole
polygon, so where it changes sign the region is cut along the line $\varphi = 0$ and only the
side $\varphi > 0$ is kept (Fig.~\ref{fig:partition}). On what remains, $p$ is the product of
two affine functions, $f = \gamma K\ell$ and $\varphi$. For a triangle of area $A_T$ with
barycentric coordinates $\lambda_k$ and vertex values $f_k$, $\varphi_k$,
\begin{equation}
  \int f\,\varphi\,\lambda_k\,\dd A
  = \frac{A_T}{60}\left(F \Phi + Q + \varphi_k F + f_k \Phi + 2 f_k \varphi_k\right),
  \label{eq:product}
\end{equation}
with $F = \sum_i f_i$, $\Phi = \sum_i \varphi_i$ and $Q = \sum_i f_i \varphi_i$, which follows from
$\int \lambda_i\lambda_j\lambda_k\,\dd A = A_T\,c_{ijk}/60$ with $c_{ijk} = 6$,
$2$ or $1$ for three, two or no equal indices. The three integrals of Eq.~\eqref{eq:product}
give the force of the triangle and its moment exactly. They are also the integrals of the linear
pressure with vertex values
\begin{equation}
  \tilde p_k = \tfrac15\left(\varphi_k F + f_k \Phi + 2 f_k \varphi_k\right) - \tfrac1{20}\left(F \Phi + Q\right),
  \label{eq:equivalent}
\end{equation}
since a linear function satisfies
$\int \tilde p\,\lambda_k\,\dd A = \tfrac{A_T}{12}(\sum_i \tilde p_i + \tilde p_k)$. Each triangle
stores $\tilde p$. The normal force and its moment are thus exact and do not depend on how a
polygon is triangulated; the friction sees the linear equivalent $\tilde p$, which is the exact
pressure whenever $\ell$ or $\varphi$ is constant on the triangle.

\subsection{Normal wrench and the side faces}

The pressure acts along the normal of the face it is on: a triangle of $\partial\region_A$ with
outward normal $\m$ pushes $A$ along $-\m$, one of $\partial\region_B$ along $+\m$. The normal
wrench on $A$ is the sum over the triangles of $\int \tilde p\,\dd A$ times that direction,
applied at $\int \tilde p\,\x\,\dd A / \int \tilde p\,\dd A$. For a flat contact the faces of
$A$ and of $B$ in $\partial\region$ are both perpendicular to $\n$, each carries half, and the
resultant is $K V \varphi(\vect{c})$ at the centre of pressure. On the thin side faces of an
overlap, which are nearly parallel to $\n$, the columns are short and the pressure has no
component along $\n$; it is a small lateral pressure of order $\delta/L$ of the normal load for
a contact of depth $\delta$ and width $L$, and it carries friction, so that the ratio of
friction to normal load is $\mu\,(1 + O(\delta/L))$: a systematic bias against pure Coulomb
friction of $10^{-3}$ to $10^{-2}$ at the depths of the battery, below the thresholds at which
the results are read. This is the difference between a pressure on the faces, which is what the
law applies, and a traction along $\n$ per unit of projected area, which is what a strict
Winkler bed would apply; the two agree to first order in the depth.

\subsection{Friction: special cases and what is exact}

Two cases fall outside the instantaneous centre of Eq.~6 of the paper. When the triangle does
not rotate in its plane ($\omega_m = 0$) the slip is uniform, and the friction is
$\mu\int p\,\dd A$ against it, applied at the centre of pressure of the triangle. When the
centre is far, the three signed terms are large and nearly cancel, so the same uniform-slip
formula is used whenever $|\omega_m|\sqrt{2A_T} \le 10^{-3}|\vv_{t0}|$, with the direction of
the slip taken at the centre of pressure. At the switch the force changes by a relative
$10^{-6}$ and the moment of one triangle jumps by up to about $3\times10^{-4}$ of
$\mu\int p\,\dd A$ times its size: far below every other error of the paper, but a
discontinuity.

The friction is exact for the pressure $\tilde p$ of each flat triangle. It is not exact where
$\tilde p$ is not the true pressure, that is, where both $\ell$ and $\varphi$ vary over a
triangle; on a faceted body spinning in a faceted seat this leaves a residual torque of the
nature of a rolling resistance, of order $h\,\omega\,b^2/(12\,\mu\,r)$ of the friction torque
for facets of width $b$ at a radius $r$, which decreases with the square of the facet width. A
face that is a seam only in part counts as surface.

The check of the closed form (\code{paper/checks/check\_polar\_moments.py}) re-implements the
formulas of the paper's Section~3.4 and Appendix~A as printed and compares them with
quadrature that does not use them, for $\q$ outside, inside, on an edge and at a vertex of the
triangle; the moments agree to $10^{-14}$ with a Gauss--Legendre rule where the integrand is
analytic and to the accuracy of a midpoint rule ($10^{-5}$) elsewhere.
% The numbers of this paragraph are those the script prints; they are not battery results.

\subsection{Gonthier's law: stick--slip and the variant of the 2018 thesis}

The force $\Fn = K V(1 - h\,v_n(\vect{c}))$ is exact volumetric damping, because $v_n$ is an
affine function of position and its integral over $\region$ is $V$ times its value at the
centroid. $\Fn$ is clamped at zero, and the rolling torque is not applied while the clamp is
active, since the pressure whose moment it is has vanished. Between pure sliding and pure
spinning the friction satisfies
$(|\vect{f}_t|/\mu\Fn)^2 + (|\vect{M}_s|/\mu\Fn r_g)^2 = 1$.

Stick--slip follows the thesis \cite{Gonthier2007thesis} in its structure: a sticking function
$s = \exp(-v_\mathrm{avg}^2/v_s^2)$ blends a bristle force and a torsional bristle torque (the
latter with stiffness and damping multiplied by $r_g^2$) with the sliding values. As in the
thesis, the sliding values are shares of one friction budget: the force takes the fraction
$C_{v,s} = s + (1-s)\,C_v$ of it and the spinning torque $C_{\omega,s} = s + (1-s)\,C_\omega$,
with $C_v = |\vv_t|/\max(v_\mathrm{avg}, v_\varepsilon)$,
$C_\omega = r_g|\omega_n|/\max(v_\mathrm{avg}, v_\varepsilon)$ and the floor
$v_\varepsilon = v_s/100$ that avoids the division by zero at rest; the sliding targets are
$\mu_d F_n C_{v,s}$ and $\mu_d r_g F_n C_{\omega,s}$.

The implementation departs from the thesis in four points, none of which changes the sliding or
the spinning law: the bristle is a spring--damper on the deflection blended with the sliding
force by $s$, the regularization of Dopico et al.~\cite{Dopico2011}, instead of a bristle state
integrated with its own differential equation; the dwell-time dependence of the static friction
is not implemented; the normal is the vector-area normal; and $V$ and $\tens{J}$ are exact for
the meshes, where the thesis approximates them with voxels.

The law \law{gonthier-thesis} keeps the friction of the second author's thesis
\cite{Vilela2018}. Its sliding force ignores the spin, its spinning torque is
$\vect{M}_s = -\mu_\mathrm{spin}\,(\Fn/V)\,\tens{J}\,\omega_n\n$, which is the expression of
the 2005 paper \cite{Gonthier2005} read with a constant coefficient and is therefore viscous
($\mu_\mathrm{spin}$ carries s/m), and its bristle is a spring between the centroid of the
overlap and an anchor fixed in $B$. The centroid of the overlap is not a material point: when a
body on a flat contact of width $L$ and depth $\delta$ rocks by an angle $\theta$, the centroid
moves by $L^2\theta/(12\,\delta)$: 0.22\,mm per microradian for $L = 0.4$\,m and
$\delta = 60$\,\textmu m, values of the battery, so 2\,mm for a rocking of 10\,\textmu rad. The
anchored bristle therefore stretches with nothing sliding. Gonthier's own bristle integrates
the relative velocity at the contact, a material quantity, and that is what \law{gonthier}
does.

\subsection{Stiction}

\paragraph{Transition.} The speed $u$ that drives the weight $s = \exp(-u^2/v_s^2)$ is
Gonthier's $v_\mathrm{avg}$ in \law{gonthier}; in \law{distributed} it is its counterpart for
the pressure distribution, $u^2 = |\vv_t(\vect{c}_p)|^2 + \rho^2\omega_n^2$, with $\vect{c}_p$
the centre of the normal load and $\rho$ the radius of gyration of the pressure about it.

\paragraph{Deflection.} The deflection is the slip of the material point of $A$ that is at the
friction point, relative to $B$, accumulated from one evaluation to the next and kept in the
frame of $B$; the torsional deflection accumulates the rotation of $A$ relative to $B$ about
$\n$. The friction point is the centroid of the overlap in \law{gonthier} and the centre of the
normal load in \law{distributed}. Because what is accumulated is the motion of material points
between two evaluations, rolling without slipping loads no bristle, a displacement of the
friction point caused by a change of shape of the overlap loads none either, the state is
objective, and the accumulated slip does not depend on the time step. The state of a pair is
the deflection, the twist, the relative pose at the last evaluation and the last pair normal.

\paragraph{Stiffness proportional to the load.} With $k = \mu_s|\Fn|/d$ and $c = k\tau$ (the
torsional bristle has $k\rho^2$ and $c\rho^2$, $r_g^2$ in \law{gonthier}), every contact
reaches its static limit after the same displacement whatever its load, as in Gonthier's
bristle, whose stiffness is $\mu_s/\delta_z$ times the normal force \cite{Gonthier2007thesis}.
For a body resting on its contact under its own weight the bristle has $\omega^2 = \mu_s g/d$
whatever the mass, and $\tau$ gives it a damping ratio $\tau\omega/2 \approx 0.55$ for
$\mu_s = 0.5$. The sensitivity of the results to $d$ and $\tau$ was not swept.

\paragraph{Limits.} The spring force never exceeds the static limit; when it would, the
deflection is scaled back onto the limit, which is the slip of the bristle. The damper alone
would reach the static limit at the creep velocity $d/\tau = 20$\,mm/s and brake a body that
must slide like a viscous drag, so the total stick force, spring plus damper, is capped at the
limit as well. In \law{gonthier} the limits of the bristle force and of the torsional bristle are
$\mu_s\Fn$ and $\mu_s\Fn r_g$ times the shares of Gonthier's thesis, which both tend to 1 at
stick. The static limit of \law{gonthier}, as published and as coded, is therefore a box: a
patch at rest can hold a force of $\mu_s\Fn$ and a torque of $\mu_s\Fn r_g$ at once, up to
$\sqrt2$ times what its friction ellipse allows. In \law{distributed} the translational and the
torsional bristle share the friction through the limit surface of the patch, approximated by
the ellipse of Howe and Cutkosky \cite{Howe1996},
\begin{equation}
  \left(\frac{|\vect{f}_\mathrm{stick}|}{\mu_s\Fn}\right)^2
  + \left(\frac{|M_\mathrm{stick}|}{T_\mathrm{max}}\right)^2 \le 1 ,
  \qquad T_\mathrm{max} = \mu_s\int p\,|\x - \vect{c}_p|\,\dd A ,
  \label{eq:ellipse}
\end{equation}
where $T_\mathrm{max}$, the static spinning capacity of the pressure distribution, is the
pressure-weighted polar moment $\sum(p_q D + \nabla p\cdot\vect{R})$ over the triangles, with
the centre at $\vect{c}_p$. Limiting each bristle on its own lets a patch cut into pieces, each
of which sees a twist of the body as slip plus twist, hold more than the whole patch. The stick
acts in the plane of the pair normal, at the centre of the normal load and with the capacity of
the normal load, so that it adds up over the pieces of a body. The state remains one per pair
of pieces; a single state per pair of bodies would not do, because one body pair can touch on
several planes whose normal loads cancel in the sum and define no tangent plane.

\subsection{Properties of the distributed law}

None of the following is a result of the battery. Those marked as not asserted are not covered
by a test of the library.
\begin{itemize}
\item \emph{Reduction to Gonthier's law.} For a flat contact the normal force is
  $K V(1 - h\,v_n(\vect{c}))$ at the centre of pressure, with Gonthier's rolling resistance.
\item \emph{Closed cycle and fast paths.} The elastic force of a sphere in a faceted bowl does
  no work over a closed cycle, and the fast paths give the same overlap and forces as the plain
  ones.
\item \emph{Symmetry (not asserted by a test).} Exchanging $A$ and $B$ gives the opposite wrench
  exactly for single convex pieces and only approximately for compounds, because the pair
  normal is taken from the non-seam faces of $B$.
\item \emph{Independence of the triangulation} for the normal wrench (exact partition and exact
  clamp), and \emph{independence of the decomposition}: a box sliding or spinning on a slab
  receives the same wrench on a slab cut into pieces as on the whole slab, to a relative
  difference below $10^{-9}$. The scenarios \scn{pieces} test the same property dynamically.
\item \emph{Continuity in the pose (by construction, not asserted by a test).} The surface of
  the overlap, the column lengths and the partition vary continuously with the pose, so the
  wrench does; no face switches on or off as a sample point crosses a surface.
\item \emph{Force at deep overlaps.} A body enclosed in another, or through it, keeps a force
  $K V$ along the last normal.
\end{itemize}
The elastic force is conservative while the normal of a pair is fixed and is not exactly so
where that normal turns with the overlap (a relative work of $10^{-4}$ over a closed cycle at
the rim of a faceted bowl; Elandt et al.~\cite{Elandt2019} note the same of elastic-foundation
models in general). The differentiability of the wrench was not analysed.

\subsection{Baselines: departures from the publications}

The baselines live in the same library, run on every engine like the two laws, receive the same
$K$, $h$ and $\mu$, and are subject to the same step bounds (Section~\ref{sup:steps}).

\paragraph{Friction quadratures.} The two variants of \law{distributed} change only the
friction: the geometry, the pressure and its clamp, the normal wrench and the stiction are the
same code. The vertex rule evaluates the friction $-\mu\,p_i\,\vv_{t,i}/|\vv_{t,i}|$ at the
three vertices of each triangle (zero at a vertex that does not slip) and interpolates it
linearly, which gives the force $\tfrac{A_T}{3}\sum_i\vect{f}_i$ and the moment of a linear
field. The equations of Hasegawa and Sato \cite{Hasegawa2004} are used as we understood them
from a reading of the paper made for the review of the literature; the paper was not accessible
to us again, so the rule was not re-checked against its Eqs.~13 to~16, and it should be read as
the vertex interpolation that the paper describes and not as a verified transcription. Their
model, which projects the triangles on one contact plane, is not reproduced, only its friction
quadrature.

\paragraph{The per-face foundation.} In \law{pcm} both bodies carry springs, each with half the
stiffness, so that a flat contact carries $K V$ as the other laws do; the force of a spring is
never negative, and the friction is a continuous function of the slip velocity with the
transition velocity $v_s$ of the other laws. The law has no stiction and keeps no state, as its
sources have none (Simbody's documentation says that the model cannot represent stiction): a
body loaded below its friction limit creeps. The implementation departs from the publications
\cite{Hippmann2004,Hippmann2024,Sherman2011,SimbodyEFF} in four points, and the first two
favour the baseline. (i) The elements are the triangles of the exact contact surface, the part
of each face inside the other body, and not the whole faces whose centroid is inside: with the
two triangles per face of the battery's grounds the published sampling would leave them without
contact, and the publications ask for a fine mesh instead. What remains of the published
discretization is one sample of depth, velocity and slip per element. (ii) The penetration is
measured to the nearest point of the other surface, as in Simbody, and not along the normal of
the element. (iii) The damping is the product form $\tfrac12 K A_T\,x\,(1 + h\,\dot x)$; the
polygonal contact model has a constant areal damping. (iv) There is no viscous friction term.

\subsection{Cost of a pair and of the laws}

The cost of one pair is that of clipping the faces of each shape that lie in the candidate
region against the other shape; it grows with the number of faces in the region and not with the
size of the meshes, and it does not depend on the depth of the overlap. Pairs are independent
and are evaluated in parallel, with one scratch workspace per thread so that no memory is
allocated after the first step, and their wrenches are accumulated in a fixed order so that the
result does not depend on the number of threads.

The geometry is shared by the two laws. Gonthier's law adds a constant amount of arithmetic per
pair. The distributed law adds, per polygon of the contact surface, the search for the faces
where its columns can end, the partition, and per resulting triangle the closed forms of
the product integral and of the polar moments; its cost grows with the number of faces in
contact. Curved bodies meshed with many small faces are the expensive case: the cap of a
64-segment sphere on a plane has tens of triangles in contact, a box on a plane has two. The
measured costs are in Section~\ref{sup:cost}.
